\chapter{能量：形式与守恒}\label{ch:g10:energy-conservation}

行驶的汽车、带电的手机、热的烤箱和举起的锤子，
有一个共同点：它们都能让某件事发生。物理学用一个数
来度量这种能力——能量——并依照一条铁律记账：总量从不改变，
它只会转移，并换一种形式。本章将打开这本账：
能量所呈现的各种形式、它流经的能量转化链，
以及一千瓦时能做些什么。

\section{能量：通用货币}

\begin{definition}[能量与焦耳]\label{def:g10:energy-conservation:energy}
\emph{能量}\index{能量}度量系统产生变化的
能力：使物质运动、将其举高、加热或使其发光。
能量的单位是\emph{焦耳}\index{焦耳}（\unit{J}）。能量可以储存、传递、转化——
物理学的每个分支都用这同一种货币来计量。
\end{definition}

\begin{definition}[动能]\label{def:g10:energy-conservation:kinetic}
质量为 $m$（单位为 \unit{kg}）的物体，以速率 $v$（单位为
\unit{m/s}）运动时具有\emph{动能}\index{动能}
$E_k = \tfrac12\, m v^2$：质量加倍，动能加倍；
速率加倍，动能则变为四倍。
\end{definition}

\begin{example}[汽车与行人]\label{ex:g10:energy-conservation:kinetic}
一辆 \qty{1300}{kg} 的汽车以 \qty{100}{km/h} $= \qty{27.8}{m/s}$ 行驶，其动能为
$E_k = \tfrac12 \times 1300 \times 27.8^2 \approx \qty{5.0e5}{J}$；
一名 \qty{70}{kg} 的行人以 \qty{1.5}{m/s} 行走，其动能为
$\tfrac12 \times 70 \times 1.5^2 \approx \qty{79}{J}$——
相差约六千倍。速率的代价很高：平方关系决定了这一点。
\end{example}

\begin{definition}[重力势能]\label{def:g10:energy-conservation:potential}
质量为 $m$ 的物体位于选定参考平面上方 $h$（单位为 \unit{m}）处，
具有\emph{重力势能}\index{势能!重力}
$E_p = m g h$，
其中 $g = \qty{9.81}{N/kg}$——该式适用于地面附近，
此时 $g$ 可视为常量。只有高度\emph{之差}才有意义：
$E_p = 0$ 的高度可以任意选取。
\end{definition}

\begin{example}[水塔]\label{ex:g10:energy-conservation:potential}
每立方米（\qty{1000}{kg}）的水储存在高 \qty{30}{m} 的水塔中，
其能量为 $1000 \times 9.81 \times 30 \approx \qty{2.9e5}{J}$，
——约为上例中汽车动能的十分之六。城市
在高处储存能量。
\end{example}

\begin{definition}[其他形式]\label{def:g10:energy-conservation:forms}
能量还隐藏在一些不那么显眼的形式中：\emph{热能}\index{热能}，
即物体分子无序的热运动（温度越高，热运动越剧烈）；
\emph{化学能}\index{化学能}，储存于分子键中，
存在于食物、木材、汽油和电池中；
\emph{电能}\index{电能}，由电流
从发电厂输送到插座；
\emph{辐射能}\index{辐射能}，即使穿过真空，也能由光携带——
太阳的能量正是这样到达我们身边的；
\emph{核能}\index{核能}，储存在原子核内部，
在恒星中释放（见\cref{ch:g12:nuclear-energy}）。
\end{definition}

\section{能量传递与能量转化链}

\begin{definition}[能量传递、能量转化与能量转化链]\label{def:g10:energy-conservation:transfer}
\emph{能量传递}\index{能量传递}是指能量在系统之间移动——
即在\emph{能量库}\index{能量库}之间移动；
\emph{能量转化}\index{能量转化}则是指能量的形式发生改变。
进行上述任一过程的装置（涡轮、电动机、灯泡、肌肉）
称为\emph{转换器}\index{转换器}。
\emph{能量转化链}\index{能量转化链}描绘了这一过程：
方框表示能量库和形式，带标签的箭头表示能量传递。
\end{definition}

\begin{omfigure}
\begin{tikzpicture}[font=\small,
  box/.style={draw, rounded corners, fill=omDef!10, minimum height=10mm, align=center},
  lab/.style={font=\footnotesize, align=center}]
  \node[box] (res)  at (0,0)   {能量库（水）\\(势能)};
  \node[box] (tur)  at (4.1,0) {转动的涡轮\\(动能)};
  \node[box] (grid) at (8.2,0) {电网中的住户\\(电能)};
  \draw[-Stealth, thick, omThm] (res) -- (tur)
    node[midway, above, lab] {水沿压力管道下落};
  \draw[-Stealth, thick, omThm] (tur) -- (grid)
    node[midway, above, lab] {交流发电机};
  \draw[-Stealth, omExo] (tur.south) -- ++(0,-0.7)
    node[below, lab] {热量（摩擦、\\湍流）};
  \draw[-Stealth, omExo] (grid.south) -- ++(0,-0.7)
    node[below, lab] {导线中的热量};
\end{tikzpicture}

{\small 水电站的能量转化链：方框是能量库，
箭头是能量传递——每一条真实的箭头都会泄漏热量。}
\end{omfigure}

\begin{method}[绘制能量转化链]\label{met:g10:energy-conservation:chain}
\begin{enumerate}
  \item 确定初始能量库及其形式（消耗的是什么？）。
  \item 沿能量转换器逐一追踪：每次传递画一个箭头，
        并在每个方框中注明形式。
  \item 在每个真实的转换器处添加一个热量箭头——总有一部分能量
        以热量的形式泄漏（摩擦、导线、废气）。
\end{enumerate}
\end{method}

单摆是最简单的能量转化链——摆至最高处时全部为势能，
经过最低点时则全部为动能，
并取该处 $E_p = 0$：

\begin{omfigure}
\begin{tikzpicture}[font=\small]
  \draw[-Stealth] (0,0) -- (0,2.5) node[above] {$E$}; \draw (0,0) -- (2.8,0);
  \draw[dashed, omExo] (0,2.0) -- (2.8,2.0) node[right] {总能量};
  \fill[omDef!60] (0.5,0) rectangle (1.1,2.0); \draw[omThm, very thick] (1.7,0) -- (2.3,0);
  \node at (0.8,-0.4) {$E_p$}; \node at (2.0,-0.4) {$E_k$};
  \node at (1.4,-1.0) {摆动最高处};
\end{tikzpicture}\qquad
\begin{tikzpicture}[font=\small]
  \draw[-Stealth] (0,0) -- (0,2.5) node[above] {$E$}; \draw (0,0) -- (2.8,0);
  \draw[dashed, omExo] (0,2.0) -- (2.8,2.0) node[right] {总能量};
  \draw[omDef, very thick] (0.5,0) -- (1.1,0); \fill[omThm!60] (1.7,0) rectangle (2.3,2.0);
  \node at (0.8,-0.4) {$E_p$}; \node at (2.0,-0.4) {$E_k$};
  \node at (1.4,-1.0) {最低点};
\end{tikzpicture}

{\small 单摆在两个位置的能量分配：$E_p$ 与 $E_k$
两者互换，其和（虚线）保持不变。}
\end{omfigure}

\section{能量守恒}

\begin{theorem}[能量守恒]\label{thm:g10:energy-conservation:conservation}
\index{能量守恒}
孤立系统 --- 即与外界不发生任何交换的系统 --- 的总能量
保持不变。能量既不会被创生，也不会
被消灭：它只改变形式与位置，一个能量库
所失去的，其余能量库恰好等量增加。
\end{theorem}
\admitted

\begin{remark}\label{rem:g10:energy-conservation:admitted}
从未有实验发现这一原理失效；本书在此将其视为给定，
并用它做数值计算。在\cref{ch:g12:work-and-energy}
中由牛顿定律推出其机械部分（$E_k + E_p$）；
完整的核算留待本科分册。
\end{remark}

\begin{example}[以焦耳衡量的自由落体]\label{ex:g10:energy-conservation:fall}
一个 \qty{2.0}{kg} 的球从 \qty{10}{m} 高处下落（空气阻力
可忽略）：开始时
$E_p = 2.0 \times 9.81 \times 10 \approx \qty{196}{J}$，落地时
$E_k = \qty{196}{J}$，于是 $v = \sqrt{2 \times 196/2.0} =
\qty{14}{m/s} \approx \qty{50}{km/h}$。不必分析受力，也无需秒表：
仅凭守恒即可算出撞击时的速率。
\end{example}

\section{效率}

\begin{definition}[效率]\label{def:g10:energy-conservation:efficiency}
转换器消耗 $E_{\text{\, 消耗}}$，并输出所需的
形式 $E_{\text{有用}}$；其 \emph{效率}\index{效率} 为
\[
\eta = \frac{E_{\text{有用}}}{E_{\text{消耗}}} < 1 .
\]
守恒禁止 $\eta > 1$，而实际转换器永远达不到 $1$：
损失的能量并未被消灭 --- 而是以热量的形式散失。
\end{definition}

\begin{example}[汽油的去向]\label{ex:g10:energy-conservation:engine}
一千克汽油储存了 \qty{45}{MJ} 的化学能；一台
汽车发动机以 $\eta \approx 0.30$ 的效率转化它：\qty{13.5}{MJ} 的
运动，\qty{31.5}{MJ} 加热了排气管、散热器和
路面。每箱油的三分之二都在使大气层变暖。
\end{example}

\begin{omfigure}
\begin{tikzpicture}[font=\small]
  \fill[omDef!40] (0,0) rectangle (3,1.5);
  \node[align=center] at (1.5,0.75) {汽油\\\qty{45}{MJ}};
  \fill[omMeth!60] (3,1.05) -- (5.6,1.55) -- (5.6,2.0) -- (3,1.5);
  \fill[omMeth!60] (5.6,1.45) -- (6.3,1.775) -- (5.6,2.1);
  \node[right] at (6.35,1.78) {运动：\qty{13.5}{MJ} ($30\%$)};
  \fill[omExo!70] (3,0) -- (5.6,-1.0) -- (5.6,0.05) -- (3,1.05);
  \fill[omExo!70] (5.6,-1.1) -- (6.3,-0.475) -- (5.6,0.15);
  \node[right] at (6.35,-0.48) {热量：\qty{31.5}{MJ} ($70\%$)};
\end{tikzpicture}

{\small 汽车发动机的能量收支，箭头宽度按比例绘制：在
\qty{45}{MJ} 的消耗中，\qty{13.5}{MJ} 推动汽车，其余均为热量。}
\end{omfigure}

\begin{remark}[效率相乘]\label{rem:g10:energy-conservation:chain-eta}
串联的转换器效率相乘：一座发电厂
（$0.38$）向电网供电（$0.94$），电网再向充电器供电（$0.85$）
的总效率为 $\eta = 0.38 \times 0.94 \times 0.85 \approx 0.30$。长
链条如同漏水的管道：每个箭头都在收走一份能量。
\end{remark}

\section{功率、瓦特与千瓦时}

\begin{definition}[功率]\label{def:g10:energy-conservation:power}
\emph{功率}\index{功率}是能量流动的快慢，
$P = E/t$，单位为 \emph{瓦特}\index{瓦特}（\unit{W}），且
$\qty{1}{W} = \qty{1}{J/s}$。能量是一个量，功率则是一种流动 ---
好比浴缸与水龙头之别。
\end{definition}

\begin{example}[水壶与人]\label{ex:g10:energy-conservation:power}
一把 \qty{2200}{W} 的水壶工作 \qty{150}{s}，
$E = 2200 \times 150 = \qty{3.3e5}{J}$。人每天大约靠
\qty{10}{MJ} 的食物维持，折合平均
功率 $\num{1.0e7}/\num{86400} \approx \qty{116}{W}$：你
就像一只明亮的旧灯泡，一直空转着。
\end{example}

\begin{definition}[千瓦时]\label{def:g10:energy-conservation:kwh}
\emph{千瓦时}\index{千瓦时}（\unit{kW.h}）是
以 \qty{1}{kW} 的功率持续运行一小时所输送的能量：
$\qty{1}{kW.h} = 1000 \times 3600\,\unit{J} = \qty{3.6}{MJ}$ ---
即电表所计的单位：功率（\unit{kW}）乘以小时数。
\end{definition}

\begin{example}[读懂电费账单]\label{ex:g10:energy-conservation:bill}
账单就是一次减法和一次乘法：新的电表读数
减去旧读数，得千瓦时数；乘以单价（约
$0.25$ 欧元每 \unit{kW.h}），再加上固定月租，即得
总额。一台 \qty{2}{kW} 的取暖器每晚运行 \qty{2.5}{h}，
每天耗电 \qty{5}{kW.h}：仅这一台每月就要 $38$ 欧元。
\end{example}

\begin{example}[数量级]\label{ex:g10:energy-conservation:orders}
值得牢记：将一个苹果举高一米，约 \qty{1}{J}；
手机电池 $\qty{10}{W.h} \approx \qty{4e4}{J}$；一碗
麦片，\qty{1.5}{MJ}；一天的食物，\qty{10}{MJ}；一千克
汽油，\qty{45}{MJ} --- 几乎恰好等于一个家庭每天的
用电量（$\approx \qty{12}{kW.h}$）；一次闪电，\qty{5}{GJ}：
相当于一百个家庭一天的用电量，一瞬间释放出来。
\end{example}

\begin{omfigure}
\begin{tikzpicture}[font=\small, y=0.5cm]
  \draw[-Stealth] (0,-0.5) -- (0,10.8) node[above] {$E$ (\unit{J})};
  \foreach \n in {0,2,4,6,8,10}
    {\draw (-0.12,\n) -- (0.12,\n); \node[left, font=\footnotesize] at (-0.15,\n) {$10^{\n}$};}
  \fill[omDef] (0,0) circle (1.6pt) node[right=4pt] {举高 \qty{1}{m} 的苹果};
  \fill[omDef] (0,4.56) circle (1.6pt) node[right=4pt] {手机电池（\qty{10}{W.h}）};
  \fill[omDef] (0,6.18) circle (1.6pt) node[right=4pt] {一碗麦片};
  \fill[omDef] (0,7) circle (1.6pt) node[right=4pt] {一天的食物（\qty{10}{MJ}）};
  \fill[omThm] (0,7.65) circle (1.6pt) node[right=4pt] {\qty{1}{kg} 汽油 $\approx$ 一个家庭一天的用电量};
  \fill[omThm] (0,9.7) circle (1.6pt) node[right=4pt] {闪电（$\approx\qty{5}{GJ}$）};
\end{tikzpicture}

{\small 能量阶梯：每一级相差 $100$ 倍。
日常生活跨越十个数量级。}
\end{omfigure}

\section{练习}

\begin{exercise}[$\star$]\label{exo:g10:energy-conservation:1}
一辆 \qty{1200}{kg} 的汽车以 \qty{90}{km/h}（即 \qty{25}{m/s}）行驶：
求其动能。将速率提高到 \qty{130}{km/h}（\qty{36.1}{m/s}）后重新计算；
速率增加 $44\%$，$E_k$ 变为原来的多少倍？
\end{exercise}

\begin{exercise}[$\star$]\label{exo:g10:energy-conservation:2}
一名质量为 \qty{72}{kg} 的登山者从海拔 \qty{1200}{m} 攀登至
\qty{1650}{m}。求其增加的势能，并将其换算为
千瓦时，并按每 \unit{kW.h} $0.25$ 欧元计算其电费。
\end{exercise}

\begin{exercise}[$\star$]\label{exo:g10:energy-conservation:3}
换算：\qty{1}{kW.h} 换算为焦耳；\qty{12}{W.h} 的手机电池换算为
千焦；\qty{45}{MJ}（一千克汽油）换算为
千瓦时；一天的 \qty{10}{MJ} 食物换算为千瓦时。
\end{exercise}

\begin{exercise}[$\star$]\label{exo:g10:energy-conservation:4}
一台 \qty{2200}{W} 的水壶工作 \qty{150}{s}。求其消耗的能量，
分别以焦耳和千瓦时表示。再求
一个人每天消耗 \qty{10}{MJ} 时的平均功率。
\end{exercise}

\begin{exercise}[$\star$]\label{exo:g10:energy-conservation:5}
写出以下各项的能量转化链（见\cref{met:g10:energy-conservation:chain}）：
自行车发电机为灯泡供电；燃气热水器；太阳能
电池板为手机充电。标出热量损失。
\end{exercise}

\begin{exercise}[$\star\star$]\label{exo:g10:energy-conservation:6}
一台发动机燃烧 \qty{2.0}{kg} 汽油（\qty{45}{MJ/kg}），
输出 \qty{27}{MJ} 的机械能。求其效率及
以热量形式释放的能量。这些热量去了哪里？
\end{exercise}

\begin{exercise}[$\star\star$]\label{exo:g10:energy-conservation:7}
一名质量为 \qty{70}{kg} 的跳崖运动员从 \qty{20}{m} 高的悬崖跳下（忽略
空气阻力）。利用能量守恒，求其入水时的动能
与入水速率，分别以 \unit{m/s} 和 \unit{km/h} 表示。
\end{exercise}

\begin{exercise}[$\star\star$]\label{exo:g10:energy-conservation:8}
3月1日电表读数为 \qty{41236}{kW.h}，30
天后为 \qty{41562}{kW.h}。求消耗的能量、日均用电量、按每
\unit{kW.h} $0.26$ 欧元加 $12$ 欧元固定费用计算的电费，以及以瓦特表示的平均功率。
\end{exercise}

\begin{exercise}[$\star\star$]\label{exo:g10:energy-conservation:9}
一个旧的 \qty{60}{W} 灯泡只将其耗电的 $5\%$ 转化为
光能；一个发光二极管在 \qty{9}{W} 时可发出相同的光。求
灯泡的光功率与发光二极管的效率；再按每天 \qty{4}{h} 计算
两者的年耗电量，以及按每
\unit{kW.h} $0.25$ 欧元计算所节省的电费。
\end{exercise}

\begin{exercise}[$\star\star$]\label{exo:g10:energy-conservation:10}
一座水坝每秒使 \qty{150}{m^3} 的水（\qty{1.5e5}{kg}）
从 \qty{40}{m} 的落差处下泄。求每秒释放的势能
——即可用功率——再求效率为
$\eta = 0.90$ 时的电功率，以及它可供给多少户平均功率为 \qty{525}{W} 的家庭。
\end{exercise}

\begin{exercise}[$\star\star$]\label{exo:g10:energy-conservation:11}
一个质量为 \qty{0.20}{kg} 的摆球从高于
最低点 \qty{5.0}{cm} 处释放：求其在最低点的速率。数分钟后它
静止悬挂：其能量去了哪里？为何这并不与
\cref{thm:g10:energy-conservation:conservation}相矛盾？
\end{exercise}

\begin{exercise}[$\star\star\star$]\label{exo:g10:energy-conservation:12}
一座燃煤电厂（$\eta_1 = 0.38$）、电网（$\eta_2 = 0.94$）和
充电器（$\eta_3 = 0.85$）为一块 \qty{12}{W.h} 的手机电池充电。
求总效率；求电厂消耗的
化学能，分别以 \unit{W.h} 和千焦表示；并求煤（\qty{25}{MJ/kg}）
的用量：每次充电所需，再求每天充电一年所需。
\end{exercise}

\begin{exercise}[$\star\star\star$]\label{exo:g10:energy-conservation:13}
一次闪电约携带 \qty{5}{GJ} 的能量。换算为
千瓦时，并换算为一户每天耗电 \qty{12.6}{kW.h} 的家庭的用电天数。
闪电在约 \qty{100}{\micro s} 内释放：求其功率，
将其与人类的平均电功率（\qty{3}{TW}）比较，
并说明为何无人收集闪电能量。
\end{exercise}

\begin{exercise}[$\star\star\star$]\label{exo:g10:energy-conservation:14}
一辆质量为 \qty{1300}{kg} 的汽车从 \qty{110}{km/h}（\qty{30.6}{m/s}）
刹车至静止。求耗散的动能；它以何种形式
呈现，又在何处？一辆电动汽车将其中 $60\%$ 回收至电池：
求每次停车回收的能量，以及需要停车多少次才能充入一千瓦时？
\end{exercise}

\begin{exercise}[$\star\star\star$]\label{exo:g10:energy-conservation:15}
骑车人与自行车（合计 \qty{85}{kg}）爬上一座山口，
海拔升高 \qty{1200}{m}；肌肉转化食物的
效率为 $\eta \approx 0.25$。
求爬坡过程的机械能、它所需的食物能量，
折合成 \qty{1.5}{MJ} 的麦片碗数，以及
储存同样能量所需汽油的质量（\qty{45}{MJ/kg}），并加以评论。
\end{exercise}

\section{问题：为房屋供电一天}

\begin{problem}[{周末问题 --- 从麦片碗到屋顶光伏板：一户普通家庭
一天的完整能量审计，以及一焦耳的价格
因出售者而异}]
\label{pb:g10:energy-conservation:1}
某个普通的星期二，一家人要借助本章知识、
走廊里的电表和电费单背面，查清电能去向。
今晚的结论是：他们一天的用电量相当于
一千克汽油、十六平方米的日照，或者五位
日夜不停蹬车的骑车人。

\medskip
\noindent\textbf{第一部分 --- 能耗审计。}
一天的用电清单：冰箱功率 \qty{100}{W}，运行 \qty{24}{h}；热水器
功率 \qty{2400}{W}，运行 \qty{2.5}{h}；烤箱功率 \qty{2000}{W}，运行
\qty{1.0}{h}；洗衣机功率 \qty{500}{W}，运行 \qty{2.0}{h}；照明
功率 \qty{60}{W}，运行 \qty{5.0}{h}；屏幕功率 \qty{150}{W}，运行 \qty{6.0}{h}。
\begin{enumerate}
  \item 分别计算各电器一天的耗电量，单位为 \unit{kW.h}。
  \item 分别以千瓦时和兆焦为单位，计算全天总耗电量；
        并与一千克汽油（\qty{45}{MJ}）相比较。
  \item 求该房屋在 \qty{24}{h} 内的平均功率；
        这是一只正在工作的水壶功率的几分之一？
  \item 哪一台电器占主导？重新计算全天总量
        若热水器仅运行 \qty{1.25}{h}。
  \item 待机指示灯长期耗电功率为 \qty{15}{W}：求其一天的耗电量、
        占全天总量的比例，以及按每 \unit{kW.h} $0.25$ 欧元计算的全年电费？
\end{enumerate}

\medskip
\noindent\textbf{第二部分 --- 电费账单。}
\begin{enumerate}[resume]
  \item 30 天内电表读数从 \qty{56214}{kW.h} 变为
        \qty{56604}{kW.h}：求用电量，
        并将日均值与审计结果相核对？
  \item 计算电费：每 \unit{kW.h} $0.25$ 欧元，再加上
        固定服务费 $14$ 欧元。
  \item 电网供电的每兆焦价格是多少？按汽油计，
        汽油每升 $1.80$ 欧元、每升储存 \qty{34}{MJ}，其每兆焦价格是多少？
  \item 一只 \qty{1.5}{MJ} 的麦片碗约 $0.50$ 欧元：其
        每兆焦价格是多少？请将三种焦耳的卖家排序。
\end{enumerate}

\medskip
\noindent\textbf{第三部分 --- 屋顶光伏。}
\begin{enumerate}[resume]
  \item 充足阳光下，每平方米约接收 \qty{1000}{W}
        的功率；光伏板以 $\eta = 0.20$ 的效率转换，求
        \qty{1}{m^2} 面积在满日照下的电功率是多少？
  \item 该地区每天平均相当于 \qty{4.0}{h} 的满日照，
        求 \qty{1}{m^2} 面积的日发电量是多少？
  \item 一家人一天 \qty{12.6}{kW.h} 的用电量需要多大的光伏板面积？
        南向屋顶可提供 \qty{20}{m^2}，能否装得下？
  \item 十二月份，等效满日照时长降至 \qty{1.5}{h}：求
        光伏板的发电量、其占需求的比例，以及由什么来填补
        不足的部分？
  \item 制造光伏板每平方米耗能约 \qty{700}{kW.h}：
        利用第 11 题的结果，求能量回收期？
\end{enumerate}

\medskip
\noindent\textbf{第四部分 --- 麦片碗。}
\begin{enumerate}[resume]
  \item 一辆发电自行车可发出 \qty{100}{W} 的电功率：
        蹬多少小时才能产生 \qty{12.6}{kW.h}——需要多少人
        连续蹬车 \qty{24}{h}？
  \item 肌肉的效率为 $\eta \approx 0.25$：求产生这些
        \qty{12.6}{kW.h} 电能所需的食物能量，以兆焦和麦片碗数表示？
  \item 蹬车一天的成本（麦片碗按每只 $0.50$ 欧元计）
        与电网对同一天的收费相比如何？
  \item 一次 \qty{5}{GJ} 的闪电相当于多少天的房屋用电量？
        已知其持续时间为 \qty{100}{\micro s}，为何它仍然
        不能为房屋供电？
  \item 焦耳从不会消失，然而电网、汽油和食物的焦耳
        售价差别很大：你实际上为之付费的是什么？
  \item 总结：用一句话概括这一天的结论——一
        千克汽油、十六平方米的晴朗屋顶，或
        五位日夜蹬车的骑车人——并指出两个
        今后切勿混淆的量（一个是存量，一个是流量）。
\end{enumerate}
\end{problem}
